\subsection{INFINITESIMAL AUTOMORPHISMS}

Lemma 1. Let G be a Lie group and \(\alpha\) a vector field on G. For all \(g \in \mathbf{G}\) , let

\[
\beta (g) = \alpha (g) g ^ {- 1} \in \mathrm{L} (\mathrm{G}).
\]

The following conditions are equivalent:

\begin{enumerate}
\def\labelenumi{(\roman{enumi})}
\item
  \(\alpha\) is a homomorphism of the group G into the group T(G);
\item
  for all \(g, g'\) in \(\mathbf{G}\) , \(\alpha(gg') = \alpha(g)g' + g\alpha(g')\) ;
\item
  for all \(g, g'\) in \(\mathbf{G}\) , \(\beta(gg') = \beta(g) + (\mathrm{Ad}g)\beta(g')\) .
\end{enumerate}

Condition (i) means that, for all \(g, g'\) in \(G\) , we have in the group \(T(G)\) :

\[
\beta (g) g \beta (g ^ {\prime}) g ^ {\prime} = \beta (g g ^ {\prime}) g g ^ {\prime}
\]

or

\[
\beta (g) ((\mathrm{Ad} g) \beta (g ^ {\prime})) g g ^ {\prime} = \beta (g g ^ {\prime}) g g ^ {\prime}.
\]

But the product of \(\beta(g)\) and (Ad \(g)\beta(g')\) in T(G) is just the sum of \(\beta(g)\) and (Ad \(g)\beta(g')\) in L(G) (§ 2, no. 1, Proposition 2). Hence (i) \(\Leftrightarrow\) (iii). On the other hand, condition (ii) may be written \(\beta(gg')gg' = \beta(g)gg' + g\beta(g')g'\) , or

\[
\beta (g g ^ {\prime}) = \beta (g) + (\mathrm{Ad} g) \beta (g ^ {\prime})
\]

and hence (ii) \(\Leftrightarrow\) (iii).

DEFINITION 1. Let G be a Lie group. An infinitesimal automorphism of G is any analytic vector field on G satisfying the conditions of Lemma 1.

Lemma 2. Let \(\mathbf{K}'\) be a non-discrete closed subfield of \(\mathbf{K}\) , A a \(\mathbf{K}'\) -manifold, B and C K-manifolds and \(f\) a \(\mathbf{K}'\) -analytic mapping of \(\mathbf{A} \times \mathbf{B}\) into C. Suppose that, for all \(a \in \mathbf{A}\) , the mapping \(b \mapsto f(a, b)\) of \(\mathbf{B}\) into \(\mathbf{C}\) is K-analytic. Then, for all \(t \in \mathbf{TA}\) , the mapping \(u \mapsto (\mathrm{Tf})(t, u)\) of TB into TC is K-analytic.

We fix \(t \in \mathrm{TA}\) and write \(g(u) = (\mathrm{Tf})(t,u)\) . Clearly \(g\) is K'-analytic. By Differentiable and Analytic Manifolds, R, 5.14.6, it suffices to prove that the tangent mappings to \(g\) are K-linear. It can be assumed that A, B, C are open neighbourhoods of 0 in complete normable spaces E, F, G over K', K, K and that \(t\) is tangent to A at 0. We identify TA, TB, TC, with A × E, B × F, C × G and \(t\) with an element of E. Then for all \((x,y) \in \mathrm{TB} = \mathrm{B} \times \mathrm{F}\) .

\[
g (x, y) = \left(f (0, x), \left(\mathrm{D} _ {1} f\right) (0, x) (t) + \left(\mathrm{D} _ {2} f\right) (0, x) (y)\right).
\]

We identify \(\mathrm{T(B\times F)}\) with \((\mathbf{B}\times \mathbf{F})\times (\mathbf{F}\times \mathbf{F})\) and \(\mathrm{T(C\times G)}\) with \((\mathbf{C}\times \mathbf{G})\times (\mathbf{G}\times \mathbf{G})\) . Then, for all

\[
((x, y), (h, k)) \in \mathrm{T} (\mathrm{B} \times \mathrm{F}) = (\mathrm{B} \times \mathrm{F}) \times (\mathrm{F} \times \mathrm{F}),
\]

\((\mathrm{Tg})((x,y),(h,k)) = ((a,b),(c,d)),\) where

\[
\begin{array}{r l} & {a = f (0, x),} \\ & {b = (\mathrm{D} _ {1} f) (0, x) (t) + (\mathrm{D} _ {2} f) (0, x) (y), \qquad c = (\mathrm{D} _ {2} f) (0, x) (h),} \\ & {d = (\mathrm{D} _ {2} \mathrm{D} _ {1} f) (0, x) (t, h) + (\mathrm{D} _ {2} \mathrm{D} _ {2} f) (0, x) (y, h) + (\mathrm{D} _ {2} f) (0, x) (k).} \end{array}
\]

We now fix \((x,y)\in \mathbf{B}\times \mathbf{F}\) . We need to prove that the mapping \((h,k)\mapsto (c,d)\) of \(\mathbf{F}\times \mathbf{F}\) into \(\mathbf{G}\times \mathbf{G}\) is K-linear. As the mapping \(x\mapsto f(0,x)\) of B into C is K-analytic, the mappings

\[
\begin{array}{c} (h, k) \mapsto (\mathrm{D} _ {2} f) (0, x) (h), \qquad (h, k) \mapsto (\mathrm{D} _ {2} \mathrm{D} _ {2} f) (0, x) (y, h), \\ (h, k) \mapsto (\mathrm{D} _ {2} f) (0, x) (k) \end{array}
\]

are K-linear. On the other hand

\[
\left(\mathrm{D} _ {2} \mathrm{D} _ {1} f\right) (0, x) (t, h) = \lim _ {\lambda \in K ^ {*}, \lambda \rightarrow 0} \lambda^ {- 1} \left(\left(\mathrm{D} _ {2} f\right) (\lambda t, x) (h) - \left(\mathrm{D} _ {2} f\right) (0, x) (h)\right)
\]

and, for fixed \(\lambda\) , the mapping \(x \mapsto f(\lambda t, x)\) is K-analytic, so that the mapping \(h \mapsto (\mathrm{D}_2f)(\lambda t, x)(h)\) is K-linear.

PROPOSITION 1. Let \(\mathbf{K}'\) be a non-discrete closed subfield of \(\mathbf{K}\) , \(\mathbf{G}\) a Lie group over \(\mathbf{K}\) , \(\mathbf{V}\) a manifold over \(\mathbf{K}'\) and \((v, g) \mapsto vg\) a \(\mathbf{K}'\) -analytic mapping of \(\mathbf{V} \times \mathbf{G}\) into \(\mathbf{G}\) . Suppose that, for all \(v \in \mathbf{V}\) , the mapping \(g \mapsto vg\) of \(\mathbf{G}\) into \(\mathbf{G}\) is an automorphism of \(\mathbf{G}\) . Let \(\varepsilon\) be an element of \(\mathbf{V}\) such that \(\varepsilon g = g\) for all \(g \in \mathbf{G}\) and \(a \in T_{\varepsilon}(\mathbf{V})\) . Then the vector field \(g \mapsto ag\) on \(\mathbf{G}\) is an infinitesimal automorphism of \(\mathbf{G}\) .

For \(v \in \mathrm{V}\) , \(g_1 \in \mathrm{G}\) , \(g_2 \in \mathrm{G}\) , \(v(g_1g_2) = (vg_1)(vg_2)\) . Hence, for \(u_1 \in \mathrm{TG}\) , \(u_2 \in \mathrm{TG}\) , \(a(u_1u_2) = (au_1)(au_2)\) (§ 2, no. 1, Proposition 3). In particular, the mapping \(g \mapsto ag\) of \(G\) into \(TG\) is a group homomorphism. On the other hand, this mapping is analytic by Lemma 2.

PROPOSITION 2. Let G be a real or complex Lie group and \(\alpha\) an infinitesimal automorphism of G. There exists a law of analytic operation \((\lambda, g) \mapsto \phi_{\lambda}(g)\) of K on G with the following properties:

\begin{enumerate}
\def\labelenumi{(\arabic{enumi})}
\item
  if D is the associated law of infinitesimal operation, then D(1) = α;
\item
  for all \(\lambda \in \mathbf{K}\) , \(\phi_{\lambda} \in \operatorname{Aut} G\) .
\end{enumerate}

\begin{enumerate}
\def\labelenumi{(\alph{enumi})}
\item
  For all \(\mu > 0\) , let \(\mathbf{K}_{\mu}\) be the open ball of centre 0 and radius \(\mu\) in \(\mathbf{K}\) . For all \(g \in G\) , let \(\mathcal{F}_g\) be the set of analytic integral curves \(f\) of \(\alpha\) defined in a ball \(\mathbf{K}_{\mu}\) and such that \(f(0) = g\) . By Differentiable and Analytic Manifolds, R, 9.1.3 and 9.1.5, \(\mathcal{F}_g\) is non-empty and two elements of \(\mathcal{F}_g\) coincide on the intersection of their domains of definition; let \(\mu(g)\) be the least upper bound of the numbers \(\mu\) such that there exists an element of \(\mathcal{F}_g\) defined in \(\mathbf{K}_{\mu}\) ; there exists a unique element of \(\mathcal{F}_g\) defined in \(\mathbf{K}_{\mu(0)}\) ; we denote it by \(f_g\) .
\item
  Let \(g_1, g_2\) be in \(\mathbf{G}\) , \(f_1 \in \mathcal{F}_{g_1}\) , \(f_2 \in \mathcal{F}_{g_2}\) with \(f_1\) and \(f_2\) defined on the same ball \(\mathbf{K}_\mu\) . Then \(f_1f_2: \mathbf{K}_\mu \to \mathbf{G}\) is analytic and \((f_1f_2)(0) = g_1g_2\) . On the other hand, for all \(\lambda \in \mathbf{K}_\mu\) ,
\end{enumerate}

\[
\begin{array}{r l} (\mathrm{T} _ {\lambda} (f _ {1} f _ {2})) 1 & = (\mathrm{T} _ {\lambda} f _ {1}) 1 \cdot f _ {2} (\lambda) + f _ {1} (\lambda) \cdot (\mathrm{T} _ {\lambda} f _ {2}) 1 \quad (\S 2, \text { Proposition   7 }) \\ & = \alpha (f _ {1} (\lambda)) f _ {2} (\lambda) + f _ {1} (\lambda) \alpha (f _ {2} (\lambda)) \\ & = \alpha ((f _ {1} f _ {2}) (\lambda)) \quad \text {(Lemma   1)} \end{array}
\]

and hence \(f_{1}f_{2}\in \mathcal{F}_{g_{1}g_{2}}\) . This proves that \(\mu (g_1g_2)\geqslant \inf (\mu (g_1),\mu (g_2))\)

\begin{enumerate}
\def\labelenumi{(\alph{enumi})}
\setcounter{enumi}{2}
\tightlist
\item
  By Differentiable and Analytic Manifolds, R, 9.1.4 and 9.1.5, there exists a neighbourhood V of \(e\) in G such that \(\sigma = \inf_{g \in \mathbf{V}} \mu(g) > 0\) . Let \(h \in \mathbf{G}\) and C be its connected component. For all \(h' \in \mathbf{C}\) , \(\mu(h') \geqslant \inf(\sigma, \mu(h)) > 0\) by (b). On the other hand, the functions \(f_h'\) , where \(h' \in \mathbf{C}\) , take their values in C. By Differentiable and Analytic Manifolds, R, 9.1.4 and 9.1.5, \(\mu = +\infty\) in C and finally \(\mu = +\infty\) in G. Then let \(f_g(\lambda) = \phi_{\lambda}(g)\) for all \(g \in \mathbf{G}\) and all \(\lambda \in \mathbf{K}\) . By Differentiable and Analytic Manifolds, R, 9.1.4 and 9.1.5, the mapping \((\lambda, g) \mapsto \phi_{\lambda}(g)\) is a law of analytic operation of K on G. Clearly, if D is the associated law of infinitesimal operation, D(1) = \(\alpha\) . By (b),
\end{enumerate}

\[
\phi_ {\lambda} (g _ {1} g _ {2}) = \phi_ {\lambda} (g _ {1}) \phi_ {\lambda} (g _ {2})
\]

for all \(\lambda \in \mathbf{K}\) , \(g_{1} \in \mathbf{G}\) , \(g_{2} \in \mathbf{G}\) .

PROPOSITION 3. Suppose that K is ultrametric. Let G be a compact Lie group and \(\alpha\) an infinitesimal automorphism of G. There exist an open subgroup I of K and a law of analytic operation \((\lambda, g) \mapsto \phi_{\lambda}(g)\) of I on G with the following properties:

\begin{enumerate}
\def\labelenumi{(\arabic{enumi})}
\item
  if D is the associated law of infinitesimal operation, then D(1) = α;
\item
  for all \(\lambda \in \mathbf{I}\) , \(\phi_{\lambda} \in \operatorname{Aut} G\) .
\end{enumerate}

As G is compact, there exist an open subgroup I' of K and a law of analytic operation \((\lambda, g) \mapsto \phi_{\lambda}(g)\) of I' on G with property (1) of the proposition (§ 4, no. 7, Corollary 2 to Theorem 6). We write \(\phi_{\lambda}(g) = f_{g}(\lambda)\) for \(\lambda \in I'\) and \(g \in G\) . Then, for \(g_{1}, g_{2}\) in \(G\) and \(\lambda \in I'\) ,

\[
\begin{array}{r l} (\mathrm{T} _ {\lambda} (f _ {g _ {1}} f _ {g _ {2}})) 1 & = (\mathrm{T} _ {\lambda} f _ {g _ {1}}) 1 \cdot f _ {g _ {2}} (\lambda) + f _ {g _ {1}} (\lambda) \cdot (\mathrm{T} _ {\lambda} f _ {g _ {2}}) 1 \\ & = \alpha (f _ {g _ {1}} (\lambda)) f _ {g _ {2}} (\lambda) + f _ {g _ {1}} (\lambda) \alpha (f _ {g _ {2}} (\lambda)) \\ & = \alpha (f _ {g _ {1}} (\lambda) f _ {g _ {2}} (\lambda)) \end{array}
\]

and \((f_{g_1}f_{g_2})(0) = g_1g_2 = f_{g_1g_2}(0)\) . Hence \(f_{g_1'g_2'}(\lambda) = f_{g_1'}(\lambda)f_{g_2'}(\lambda)\) for

\[
(g _ {1} ^ {\prime}, g _ {2} ^ {\prime}, \lambda)
\]

in a neighbourhood of \((g_1, g_2, 0)\) (Differentiable and Analytic Manifolds, R, 9.1.8). As G is compact, there exists an open subgroup I of I' such that \(f_{g_1 g_2}(\lambda) = f_{g_1}(\lambda) f_{g_2}(\lambda)\) for all \(g_1 \in G\) , \(g_2 \in G\) , \(\lambda \in I\) . In other words, \(\phi_\lambda \in \operatorname{Aut} G\) for \(\lambda \in I\) .

Lemma 3. Let G and G' be Lie groups and \(\phi\) a homomorphism of G into \(\operatorname{Aut}(\mathbf{G}')\) . Let \(f(g, g') = (\phi(g))(g')\) for \(g \in \mathbf{G}\) , \(g' \in \mathbf{G}'\) . Consider the following conditions:

\begin{enumerate}
\def\labelenumi{(\roman{enumi})}
\item
  \(f\) is analytic;
\item
  \(f\) is analytic in a neighbourhood of \((e_{\mathbf{G}}, e_{\mathbf{G}^{\prime}})\) ;
\item
  for all \(g' \in G'\) , the mapping \(g \mapsto f(g, g')\) is analytic.
\end{enumerate}

Then (i) \(\Leftrightarrow\) ((ii) and (iii)). If G is connected, (i) \(\Leftrightarrow\) (ii).

Clearly (i) implies (ii) and (iii). Let \(g_0 \in G, g_0' \in G'\) . For all \(g \in G, g' \in G'\) ,

\[
f \left(g g _ {0}, g ^ {\prime} g _ {0} ^ {\prime}\right) = (\phi (g) \phi (g _ {0})) \left(g ^ {\prime} g _ {0} ^ {\prime}\right) = \phi (g) \left(\phi (g _ {0}) g ^ {\prime}\right)\cdot \phi (g) \left(\phi (g _ {0}) g _ {0} ^ {\prime}\right).
\]

This proves the implication ((ii) and (iii)) \(\Rightarrow\) (i). Finally, if \(\mathbf{G}'\) is connected, \(\mathbf{G}'\) is generated by every neighbourhood of \(e_{\mathbf{G}'}\) and hence (ii) \(\Rightarrow\) (iii).

